\documentclass[11pt]{article}
\usepackage[a4paper,margin=25mm]{geometry}
\usepackage[T1]{fontenc}
\usepackage{lmodern}
\usepackage{booktabs}
\usepackage{amsmath}
\usepackage{graphicx}
\usepackage{hyperref}
\usepackage{xcolor}
\usepackage{listings}
\lstset{basicstyle=\ttfamily\small, columns=fullflexible, keepspaces=true,
  breaklines=true, frame=single, framerule=0.3pt, rulecolor=\color{gray}}
\newcommand{\code}[1]{\texttt{#1}}

\title{\code{q0corr.py}: the seed-channel correction, step by step}
\author{BGOTSCAN}
\date{2026-10-15}

\begin{document}
\maketitle

\section{Purpose}

The figures of the note \code{docs/q0\_correction.tex}: the response
correction by seed channel that \code{offline.py} applies to the
equalized positive sum $Q^{\rm eq}_0$ to obtain $Q^{\rm pos}_0$, shown
one step at a time on one run. The script reads the per-event file of
\code{offline.py} (\code{pass/RUN\_offline.h5}: $Q^{\rm eq}_0$,
$Q^{\rm pos}_0$, the seed channel, the class and the pile-up flag) and
repeats the measurement of the photopeak position per seed channel that
\code{offline.py} makes, before and after the correction. The definitions
are those of \code{offline\_doc.tex}, "Response correction by seed
channel".

\section{Usage}

\begin{lstlisting}
python3 src/py/q0corr.py RUN [--channels 24,34,48,3] [--min-events 40] [--save NAME]
\end{lstlisting}

\begin{center}
\begin{tabular}{l l p{80mm}}
\toprule
option & default & meaning \\
\midrule
\code{RUN} & & run name, e.g. \code{195K\_41V\_th20}; reads \code{pass/RUN\_offline.h5} \\
\code{--channels LIST} & 24,34,48,3 & seed channels drawn one by one in the \code{seeds} figure \\
\code{--min-events N} & 40 & events a channel needs as seed to get a position \\
\code{--save NAME} & & save the figures as \code{src/py/plots/RUN/NAME\_q0corr\_*.png} \\
\bottomrule
\end{tabular}
\end{center}

Sample: REAL events without pile-up. For every channel $k$ that is the
seed of at least $N$ of them, the photopeak position $\mu_k$ is the
highest maximum of the smoothed 60-bin histogram of $Q^{\rm eq}_0$ over
those events, between zero and the 99.5th percentile of the whole
sample; the same on $Q^{\rm pos}_0$ gives $\mu'_k$. The terminal prints
one line per channel (events, $\mu_k$ and $\mu'_k$ relative to their
medians) and the spread of each set, rms over median.

\section{Figures}

\paragraph{\code{seedmap}.} Two maps of the array: $\mu_k$ over its
median before the correction, and $\mu'_k$ over its median after.

\paragraph{\code{seeds}.} Two panels. Left, the $Q^{\rm eq}_0$ spectrum
of the events whose seed is each of the channels of \code{--channels},
with its $\mu_k$ as a dashed line and the median as a black line; right,
the same events on $Q^{\rm pos}_0$.

\paragraph{\code{spectra}.} The whole sample on $Q^{\rm eq}_0$ and on
$Q^{\rm pos}_0$, with the photopeak fit of \code{offline.py} (Gaussian on
a step) and its width.

\section{Example}

\begin{lstlisting}
python3 src/py/q0corr.py 195K_41V_th20 --save 195K_41V     # docs/figs/q0corr_195K_{seedmap,seeds,spectra}.png
\end{lstlisting}

\end{document}
